% GENERATED FILE. Edit canonical lessons, then run npm run generate:books.
% canonical-source-hash: fnv1a64:13d5b3e55d3892af
% canonical-lessons: GE-C357-gewerkschaft, GE-C357-betriebsrat, GE-C357-forderung, GE-C357-unterstuetzen, GE-C357-unterstuetzung

\chapter{Die Gewerkschaft}
\label{ch:ge-die-gewerkschaft}
\hlchaptermodality{357}

\begin{chapteropening}
\textbf{By the end of this chapter:} \emph{I can talk about unions and works councils, workers' demands and support for them.}

Complete the last lesson of chapter 357: I can talk about unions and works councils, workers' demands and support for them.
\end{chapteropening}

\section[die Gewerkschaft]{die Gewerkschaft — the trade union}

\label{lesson:GE-C357-gewerkschaft}



\begin{quote}
Before the new one: say the German for to cope with, then the German for to fail.
\end{quote}

\subsection*{You'll want to know: die Gewerkschaft}
\textbf{die Gewerkschaft} — \textquotedblleft{}the trade union\textquotedblright{}.

\begin{quote}
Die Gewerkschaft fordert mehr Lohn für die Arbeiter.
\end{quote}

\begin{quote}
\emph{The union is demanding higher pay for the workers.}
\end{quote}

Plural: \emph{die Gewerkschaften}. When talks fail, a union may call a strike: \emph{Die Gewerkschaft ruft zum Streik auf.}

\begin{cousinweb}
\emph{Gewerkschaft} joins \emph{Gewerk} \textquotedblleft{}trade, craft\textquotedblright{}, from \emph{Werk} \textquotedblleft{}work\textquotedblright{}, Old High German \emph{werc}, Proto-Germanic \emph{werką} (EtymDB), to the ending \emph{-schaft}. It is usually said to have named a mining company before it named a workers' union.
\end{cousinweb}

\begin{grammarlens}[title={What you've built}]
One more word for telling, explaining and describing at B1.
\end{grammarlens}

\subsection*{Your turn}
\begin{itemize}
\raggedright
  \item \emph{Say it:} \emph{die Gewerkschaft}
  \item \emph{Say it:} \emph{die Gewerkschaft}, once more
  \item \emph{Say it:} a sentence of your own with \emph{die Gewerkschaft}
  \item \emph{Recall:} say the German for to cope with, then the German for to fail, then say \emph{die Gewerkschaft} again
\end{itemize}

\begin{tcolorbox}[breakable,colback=teal!4,colframe=teal!35!black,title={Before you move on}]
What does die Gewerkschaft mean? (\textquotedblleft{}The trade union\textquotedblright{}.) Say it once more.
\end{tcolorbox}
\section[der Betriebsrat]{der Betriebsrat — the works council}

\label{lesson:GE-C357-betriebsrat}



\begin{quote}
Before the new one: say the German for to fail, then the German for the trade union.
\end{quote}

\subsection*{You'll want to know: der Betriebsrat}
\textbf{der Betriebsrat} — \textquotedblleft{}the works council\textquotedblright{}.

\begin{quote}
Der Betriebsrat spricht mit dem Chef über die Arbeitszeiten.
\end{quote}

\begin{quote}
\emph{The works council talks to the boss about working hours.}
\end{quote}

Plural: \emph{die Betriebsräte}. Also a member of it: \emph{Sie ist Betriebsrätin.}

\begin{cousinweb}
\emph{Betriebsrat} joins \emph{Betrieb} \textquotedblleft{}firm\textquotedblright{}, from \emph{betreiben} \textquotedblleft{}to run\textquotedblright{}, to \emph{Rat} \textquotedblleft{}council, advice\textquotedblright{}, inherited from Old High German \emph{rāt} and Proto-Germanic \emph{rēdaz} (EtymDB).
\end{cousinweb}

\begin{grammarlens}[title={What you've built}]
One more word for telling, explaining and describing at B1.
\end{grammarlens}

\subsection*{Your turn}
\begin{itemize}
\raggedright
  \item \emph{Say it:} \emph{der Betriebsrat}
  \item \emph{Say it:} \emph{der Betriebsrat}, once more
  \item \emph{Say it:} a sentence of your own with \emph{der Betriebsrat}
  \item \emph{Recall:} say the German for to fail, then the German for the trade union, then say \emph{der Betriebsrat} again
\end{itemize}

\begin{tcolorbox}[breakable,colback=teal!4,colframe=teal!35!black,title={Before you move on}]
What does der Betriebsrat mean? (\textquotedblleft{}The works council\textquotedblright{}.) Say it once more.
\end{tcolorbox}
\section[die Forderung]{die Forderung — the demand}

\label{lesson:GE-C357-forderung}



\begin{quote}
Before the new one: say the German for the trade union, then the German for the works council.
\end{quote}

\subsection*{You'll want to know: die Forderung}
\textbf{die Forderung} — \textquotedblleft{}the demand\textquotedblright{}.

\begin{quote}
Die Forderungen der Arbeiter sind klar: mehr Geld und weniger Stunden.
\end{quote}

\begin{quote}
\emph{The workers' demands are clear: more money and fewer hours.}
\end{quote}

Plural: \emph{die Forderungen}. From \emph{fordern} — to demand. A sum someone owes you is also \emph{eine Forderung}.

\begin{cousinweb}
\emph{Forderung} is the noun of \emph{fordern} \textquotedblleft{}to demand\textquotedblright{}. EtymDB has no history for the verb; it is usually traced to Old High German \emph{fordarōn} and the family of \emph{vorder} \textquotedblleft{}front\textquotedblright{}: you call for a thing to come forward.
\end{cousinweb}

\begin{grammarlens}[title={What you've built}]
One more word for telling, explaining and describing at B1.
\end{grammarlens}

\subsection*{Your turn}
\begin{itemize}
\raggedright
  \item \emph{Say it:} \emph{die Forderung}
  \item \emph{Say it:} \emph{die Forderung}, once more
  \item \emph{Say it:} a sentence of your own with \emph{die Forderung}
  \item \emph{Recall:} say the German for the trade union, then the German for the works council, then say \emph{die Forderung} again
\end{itemize}

\begin{tcolorbox}[breakable,colback=teal!4,colframe=teal!35!black,title={Before you move on}]
What does die Forderung mean? (\textquotedblleft{}The demand\textquotedblright{}.) Say it once more.
\end{tcolorbox}
\section[unterstützen]{unterstützen — to support}

\label{lesson:GE-C357-unterstuetzen}



\begin{quote}
Before the new one: say the German for the works council, then the German for the demand.
\end{quote}

\subsection*{You'll want to know: unterstützen}
\textbf{unterstützen} — \textquotedblleft{}to support\textquotedblright{}.

\begin{quote}
Meine Eltern unterstützen mich beim Studium.
\end{quote}

\begin{quote}
\emph{My parents support me during my studies.}
\end{quote}

Inseparable: \emph{sie unterstützt, sie hat unterstützt}.

\begin{cousinweb}
\emph{unterstützen} is inherited from Middle High German \emph{understüzzen} and Old High German \emph{untarstuzzen} (EtymDB), from \emph{unter} \textquotedblleft{}under\textquotedblright{} and \emph{stützen} \textquotedblleft{}to prop up\textquotedblright{}: you put a prop under someone.
\end{cousinweb}

\begin{grammarlens}[title={What you've built}]
One more word for telling, explaining and describing at B1.
\end{grammarlens}

\subsection*{Your turn}
\begin{itemize}
\raggedright
  \item \emph{Say it:} \emph{unterstützen}
  \item \emph{Say it:} \emph{unterstützen}, once more
  \item \emph{Say it:} a sentence of your own with \emph{unterstützen}
  \item \emph{Recall:} say the German for the works council, then the German for the demand, then say \emph{unterstützen} again
\end{itemize}

\begin{tcolorbox}[breakable,colback=teal!4,colframe=teal!35!black,title={Before you move on}]
What does unterstützen mean? (\textquotedblleft{}To support\textquotedblright{}.) Say it once more.
\end{tcolorbox}
\section[die Unterstützung]{die Unterstützung — the support}

\label{lesson:GE-C357-unterstuetzung}



\begin{quote}
Before the new one: say the German for the demand, then the German for to support.
\end{quote}

\subsection*{You'll want to know: die Unterstützung}
\textbf{die Unterstützung} — \textquotedblleft{}the support\textquotedblright{}.

\begin{quote}
Ohne die Unterstützung meiner Freunde hätte ich das nicht geschafft.
\end{quote}

\begin{quote}
\emph{Without my friends' support I would not have managed it.}
\end{quote}

No plural in everyday use. \emph{Mit Unterstützung von} — with the support of.

\begin{cousinweb}
\emph{Unterstützung} is the noun of \emph{unterstützen}, which is inherited from Old High German \emph{untarstuzzen} (EtymDB), from \emph{unter} \textquotedblleft{}under\textquotedblright{} and \emph{stützen} \textquotedblleft{}to prop up\textquotedblright{}.
\end{cousinweb}

\begin{grammarlens}[title={What you've built}]
One more word for telling, explaining and describing at B1.
\end{grammarlens}

\subsection*{Your turn}
\begin{itemize}
\raggedright
  \item \emph{Say it:} \emph{die Unterstützung}
  \item \emph{Say it:} \emph{die Unterstützung}, once more
  \item \emph{Say it:} a sentence of your own with \emph{die Unterstützung}
  \item \emph{Recall:} say the German for the demand, then the German for to support, then say \emph{die Unterstützung} again
\end{itemize}

\begin{tcolorbox}[breakable,colback=teal!4,colframe=teal!35!black,title={Before you move on}]
What does die Unterstützung mean? (\textquotedblleft{}The support\textquotedblright{}.) Say it once more.
\end{tcolorbox}
